% !TEX program = lualatex
% Relatório técnico — Network Security & Cloud Essentials
% Modelo abnTeX2 (ABNT NBR 14724). Compile com LuaLaTeX:  latexmk  (usa o latexmkrc)
% Trechos em vermelho [PREENCHER: ...] são para você completar.
% As caixas "EVIDÊNCIA PENDENTE" somem quando os scripts do lab geram os arquivos em evidencias/.

\documentclass[12pt,openright,oneside,a4paper,chapter=TITLE,section=TITLE,brazil]{abntex2}

\usepackage{fontspec}
\setmainfont{TeX Gyre Termes}             % equivalente livre da Times New Roman
\setmonofont{DejaVu Sans Mono}[Scale=0.82]
\usepackage{amsmath}
\usepackage{indentfirst}
\usepackage{graphicx}
\usepackage{xcolor}
\usepackage{fvextra}
\usepackage{microtype}
\usepackage[alf]{abntex2cite}

% ---------------------------------------------------------------------------
% Macros do modelo
% ---------------------------------------------------------------------------
\newcommand{\preencher}[1]{\textcolor{red}{[PREENCHER: #1]}}
\newcolumntype{L}[1]{>{\raggedright\arraybackslash}p{#1}}   % coluna alinhada à esquerda
\newcommand{\evidir}{evidencias}

% Versão compacta (padrão) x completa: compile com -pretex='\def\completo{1}' para a completa.
\newif\ifcompleto
\ifdefined\completo\completotrue\fi

\newcounter{evidencia}[chapter]
\renewcommand{\theevidencia}{\thechapter.\arabic{evidencia}}

% \evidencia{arquivo.txt}{legenda} — insere a saída gravada pelos scripts do laboratório
\newcommand{\evidencia}[2]{%
  \par\medskip\refstepcounter{evidencia}%
  \begin{center}\small\textbf{Evidência \theevidencia} --- #2\end{center}\vspace{-0.6\baselineskip}%
  \IfFileExists{\evidir/#1}{%
    \VerbatimInput[fontsize=\scriptsize,breaklines,breakanywhere,frame=single,framesep=2mm]{\evidir/#1}%
  }{%
    \noindent\fbox{\parbox{\dimexpr\linewidth-2\fboxsep-2\fboxrule}{\centering\small
      \textcolor{red}{\textbf{EVIDÊNCIA PENDENTE:} \texttt{\evidir/\detokenize{#1}}}\\
      Execute o script correspondente do laboratório (ver Anexo~\ref{anexo:roteiro}).}}%
  }%
  \par\smallskip{\centering\footnotesize Fonte: elaborado pelo autor (\texttt{\detokenize{#1}}).\par}\medskip}

% \evidenciaextra: saída automática dos scripts que repete um print manual; só na versão completa
\newcommand{\evidenciaextra}[2]{\ifcompleto\evidencia{#1}{#2}\fi}

% \printtela{arquivo}{legenda} — insere um print de tela de figuras/ (opcional)
% \evidenciaresumo: versão compacta usa evidencias/r_<arquivo> (só títulos e vereditos, ou linhas truncadas)
\newcommand{\evidenciaresumo}[2]{\ifcompleto\evidencia{#1}{#2}\else\evidencia{r_#1}{#2 (resumo)}\fi}

\newcommand{\printw}{0.62}   % largura dos prints (fração da largura do texto)
\newcommand{\printtela}[2]{%
  \begin{figure}[htbp]\centering
    \caption{#2}
    \IfFileExists{figuras/#1}{\includegraphics[width=\printw\textwidth]{figuras/#1}}{%
      \fbox{\parbox{0.8\textwidth}{\centering\textcolor{red}{PRINT PENDENTE: figuras/\detokenize{#1}}}}}
    \legend{Fonte: elaborado pelo autor.}
  \end{figure}}

% ---------------------------------------------------------------------------
% Dados do trabalho
% ---------------------------------------------------------------------------
\titulo{Segurança de Rede e Fundamentos de Nuvem:\\ endereçamento, filtragem de tráfego, IAM e armazenamento em laboratório controlado}
\autor{Eric Kinoshita \\ RM559216}
\local{}
\data{2026}
\instituicao{Network Security \& Cloud Essentials}
\tipotrabalho{Relatório técnico}
\preambulo{Relatório técnico apresentado à disciplina Network Security \& Cloud Essentials, como requisito parcial de avaliação.}
\orientador{}

% Capa conforme NBR 14724: instituição, autor, título, local e ano
\renewcommand{\imprimircapa}{%
  \begin{capa}\center
    {\ABNTEXchapterfont\large\imprimirinstituicao\par}
    \vspace{2cm}
    {\ABNTEXchapterfont\large\imprimirautor\par}
    \vfill
    {\ABNTEXchapterfont\bfseries\LARGE\imprimirtitulo}
    \vfill
    \large\imprimirlocal\par
    \large\imprimirdata
    \vspace*{1cm}
  \end{capa}}

% Folha de rosto (NBR 14724) sem a linha de orientador
\renewcommand{\imprimirfolhaderosto}{%
  \begin{folhaderosto}
    \begin{center}{\ABNTEXchapterfont\large\imprimirautor}\end{center}
    \vfill
    \begin{center}{\ABNTEXchapterfont\bfseries\Large\imprimirtitulo}\end{center}
    \vfill
    \hspace{.45\textwidth}\begin{minipage}{.5\textwidth}\SingleSpacing\imprimirpreambulo\end{minipage}
    \vfill
    \begin{center}{\large\imprimirinstituicao}\par{\large\imprimirdata}\end{center}
  \end{folhaderosto}}

\makeatletter
\hypersetup{pdftitle={\@title}, pdfauthor={\@author}, pdfsubject={Network Security & Cloud Essentials},
  colorlinks=true, linkcolor=black, citecolor=black, urlcolor=blue, bookmarksdepth=4}
\makeatother

\makeindex

\begin{document}
\selectlanguage{brazil}
\frenchspacing

% ----------------------------- Pré-textuais --------------------------------
\imprimircapa
\imprimirfolhaderosto   % obrigatória pela NBR 14724

\pdfbookmark[0]{\contentsname}{toc}
\tableofcontents*
\cleardoublepage

% ------------------------------- Textuais ----------------------------------
\textual
\chapter{Introdução}

A proteção de um ambiente computacional combina controles em camadas distintas: a
forma como a rede é endereçada e segmentada, as regras que filtram o tráfego entre os
hosts e, nos serviços de nuvem, quem pode executar quais ações sobre quais recursos. Uma
falha em qualquer dessas camadas costuma bastar para expor dados, como mostram os
recorrentes casos de \textit{buckets} de armazenamento deixados públicos por engano.

Este relatório documenta um laboratório controlado que exercita essas três camadas. Duas
máquinas virtuais Ubuntu Server 22.04, ligadas por uma rede interna isolada do VirtualBox,
fazem os papéis de \emph{servidor} e \emph{cliente}. Sobre elas foram realizadas quatro
fases: (1) endereçamento estático e divisão da rede 192.168.50.0/24 em sub-redes /26;
(2) publicação de um serviço web na porta 8080 e de SSH, com um \textit{firewall} local
(\texttt{ufw}) que libera somente a porta 8080 vinda do cliente, validado com
\texttt{nmap} antes e depois; (3) criação de um usuário com permissão mínima sobre um
\textit{bucket} compatível com S3 e a correção proposital desse \textit{bucket} de público
para privado; e (4) uma verificação final que reúne evidências de todos os controles.

Como o projeto MinIO de código aberto, sugerido pelo enunciado para simulação local, foi
arquivado e teve os binários retirados do site oficial, a Fase 3 utilizou o SILO, um
\textit{fork} comunitário mantido do MinIO que preserva a mesma API S3, a mesma API
administrativa e o mesmo cliente de linha de comando (renomeado \texttt{mcli})
\cite{silo-docs,minio-arquivado}.

Todas as etapas foram automatizadas por \textit{scripts} que registram cada comando,
sua saída e seu código de saída; esses registros compõem as evidências apresentadas no
\autoref{cap:desenvolvimento} e no Anexo~\ref{anexo:roteiro}. O documento segue a
NBR 14724 \cite{abnt14724}.

\chapter{Desenvolvimento}\label{cap:desenvolvimento}

O ambiente foi criado com Vagrant sobre o VirtualBox (\texttt{Vagrantfile} no
Anexo~\ref{anexo:roteiro}). Cada VM possui dois adaptadores: o primeiro, em NAT, serve
apenas para instalação de pacotes e para a gerência pelo comando \texttt{vagrant ssh}; o
segundo está ligado à rede interna \texttt{lab-net}, que não tem comunicação com o
hospedeiro nem com a rede local, e é a única rede do laboratório. A
\autoref{tab:topologia} resume a topologia.

\begin{table}[htbp]
  \centering
  \caption{Topologia do laboratório}\label{tab:topologia}
  \begin{tabular}{llll}
    \hline
    \textbf{VM} & \textbf{Papel} & \textbf{Endereço (Fase 1, /24)} & \textbf{Endereço (/26)} \\
    \hline
    servidor & HTTP 8080, SSH, SILO & 192.168.50.10/24 & 192.168.50.10/26 \\
    cliente  & origem dos testes, \texttt{nmap} & 192.168.50.20/24 & 192.168.50.20/26 \\
    \hline
  \end{tabular}
  \legend{Fonte: elaborado pelo autor.}
\end{table}

% ===========================================================================
\section{Endereçamento e segmentação}\label{sec:enderecamento}

\subsection{Configuração de IP estático}

O endereço estático foi definido com o Netplan \cite{netplan-docs}, gravando o arquivo
\texttt{/etc/netplan/60-lab.yaml} na interface da rede interna e aplicando-o com
\texttt{netplan apply}. As evidências a seguir mostram o arquivo e o estado resultante
das interfaces.

\evidenciaextra{fase1_ip_servidor_24.txt}{IP estático do servidor (/24)}
\evidenciaextra{fase1_ip_cliente_24.txt}{IP estático do cliente (/24)}

A comunicação foi confirmada a partir do cliente com \texttt{ping} e com uma sessão
\texttt{ssh} autenticada por chave, que executa comandos no servidor.

\evidenciaextra{fase1_conectividade_24.txt}{Conectividade cliente $\rightarrow$ servidor na rede /24}

\printtela{fase1-conectividade-24.png}{Execução manual no cliente: IP 192.168.50.20/24, \texttt{ping} sem perdas e \texttt{ssh} ao servidor}

\subsection{Cálculo das sub-redes /26}

Na rede 192.168.50.0/24 os 24 primeiros bits identificam a rede e os 8 restantes os
\textit{hosts} \cite{rfc950,rfc4632}. Passar para /26 significa emprestar 2 bits da parte
de \textit{host}, o que gera $2^2 = 4$ sub-redes de $2^6 = 64$ endereços cada. Descontados
o endereço de rede e o de \textit{broadcast}, restam $64 - 2 = 62$ \textit{hosts} úteis por
sub-rede. A nova máscara é:
\[
  /26 = \underbrace{11111111.11111111.11111111.11}_{26\ \text{bits}}000000 = 255.255.255.192
\]
O tamanho do bloco é $256 - 192 = 64$, de modo que as sub-redes começam nos múltiplos de 64
do último octeto (\autoref{tab:subredes}).

\begin{table}[htbp]
  \centering
  \caption{Divisão de 192.168.50.0/24 em sub-redes /26}\label{tab:subredes}
  \begin{tabular}{clllc}
    \hline
    \textbf{Sub-rede} & \textbf{Endereço de rede} & \textbf{Faixa de hosts} & \textbf{Broadcast} & \textbf{Hosts} \\
    \hline
    A & 192.168.50.0/26   & .1 -- .62    & 192.168.50.63  & 62 \\
    B & 192.168.50.64/26  & .65 -- .126  & 192.168.50.127 & 62 \\
    C & 192.168.50.128/26 & .129 -- .190 & 192.168.50.191 & 62 \\
    D & 192.168.50.192/26 & .193 -- .254 & 192.168.50.255 & 62 \\
    \hline
  \end{tabular}
  \legend{Fonte: elaborado pelo autor.}
\end{table}

O enunciado pede duas sub-redes /26; foram utilizadas a A (servidor .10 e cliente .20) e
a B (usada no teste de segmentação). Vale observar que dividir a /24 em exatamente duas
metades exigiria máscara /25 (126 \textit{hosts} cada); com /26, as duas sub-redes usadas
ocupam metade do espaço e as sub-redes C e D ficam reservadas para crescimento. O cálculo
foi conferido com a biblioteca \texttt{ipaddress} do Python:

\evidencia{fase1_subredes.txt}{Cálculo das sub-redes /26 conferido com \texttt{ipaddress}}

\subsection{Aplicação da máscara /26 e teste de segmentação}

As duas VMs foram reconfiguradas para /26, permanecendo na sub-rede A, e a comunicação
foi testada novamente.

\evidenciaextra{fase1_ip_servidor_26.txt}{Servidor reconfigurado para /26}
\evidenciaextra{fase1_ip_cliente_26.txt}{Cliente reconfigurado para /26}

\printtela{fase1-servidor-24-26.png}{Execução manual no servidor: IP 192.168.50.10/24 e, em seguida, /26 com a rota 192.168.50.0/26}
\printtela{fase1-cliente-26.png}{Execução manual no cliente: passagem de 192.168.50.20/24 para /26 e rota 192.168.50.0/26}
\evidenciaextra{fase1_conectividade_26.txt}{Conectividade cliente $\rightarrow$ servidor na sub-rede A (/26)}

\printtela{fase1-conectividade-26.png}{Execução manual no cliente: \texttt{ping} e \texttt{ssh} ao servidor na sub-rede A (/26)}

Para demonstrar o efeito da segmentação, o cliente foi movido temporariamente para a
sub-rede B (192.168.50.70/26). Como não existe roteador entre A e B, o sistema não tem
rota para o servidor pela interface do laboratório e o \texttt{ping} falha; em seguida o
cliente foi devolvido à sub-rede A.

\evidenciaextra{fase1_segmentacao.txt}{Cliente na sub-rede B não alcança o servidor na sub-rede A}

\printtela{fase1-segmentacao.png}{Execução manual no cliente: em 192.168.50.70/26 (sub-rede B) a rota para o servidor sai pela NAT e o \texttt{ping} pela interface do lab tem 100\% de perda}

% ===========================================================================
\section{Configuração do firewall e varredura de portas}\label{sec:firewall}

\subsection{Serviços publicados no servidor}

O servidor web foi executado como serviço do \textit{systemd} com
\texttt{python3 -m http.server 8080} \cite{python-httpserver}, com usuário dinâmico sem
privilégios (\texttt{DynamicUser=yes}); o SSH é o OpenSSH padrão do Ubuntu.

\evidenciaextra{fase2_servicos.txt}{Serviços HTTP (8080) e SSH (22) ativos no servidor}

\printtela{fase2-servicos.png}{Execução manual no servidor: \texttt{python3 -m http.server 8080} em segundo plano, SSH ativo e as portas 22 e 8080 em escuta}

\subsection{Varredura antes do firewall}

A partir do cliente foram feitas duas varreduras TCP SYN com \texttt{nmap}
\cite{lyon2009nmap}: uma nas portas de interesse com detecção de versão e outra em todas as
65.535 portas. Sem \textit{firewall}, as portas 22 e 8080 aparecem como \emph{open}.

\evidenciaextra{fase2_nmap_antes.txt}{\texttt{nmap} do cliente contra o servidor, antes do firewall}

\printtela{fase2-nmap-antes.png}{Execução manual no cliente: \texttt{nmap} antes do firewall, com 22/tcp e 8080/tcp abertas e 65.533 portas fechadas}

\subsection{Regras de filtragem}

O \texttt{ufw} \cite{ufw-docs} foi configurado com política padrão de negar todo tráfego
de entrada e uma única regra de liberação, restrita em quatro dimensões: interface do
laboratório, origem (192.168.50.20), destino (192.168.50.10) e porta/protocolo
(8080/TCP). Isso segue a recomendação de negar por padrão e liberar apenas o
estritamente necessário \cite{nist80041}. Com essa política o SSH fica bloqueado para
qualquer origem da rede do laboratório.

Duas observações de projeto:
\begin{alineas}
  \item o adaptador NAT do Vagrant não faz parte da rede do laboratório, e bloquear tudo
    nele impediria a própria gerência da VM. Por isso foi criada uma exceção apenas para a
    porta 22 vinda do hospedeiro pela NAT (10.0.2.2), registrada com comentário na regra.
    Essa exceção não é visível a partir do cliente;
  \item o enunciado da Fase 2 pode ser lido como ``SSH permitido somente a partir do
    cliente'', mas a Fase 4 exige que \emph{somente} a porta 8080 esteja acessível a partir
    do cliente. Adotou-se a leitura da Fase 4; o \textit{script} aceita
    \texttt{PERMITIR\_SSH\_CLIENTE=1} para a leitura alternativa.
\end{alineas}

\evidenciaextra{fase2_firewall.txt}{Regras do \texttt{ufw} aplicadas no servidor}

\printtela{fase2-firewall.png}{Execução manual no servidor: regras do \texttt{ufw} e \texttt{ufw status verbose} com entrada negada por padrão e somente 8080/TCP liberada para o cliente}

No \texttt{ufw status verbose}, a política de roteamento aparece como \emph{disabled}
porque o encaminhamento de pacotes (\texttt{net.ipv4.ip\_forward}) está desligado no
servidor: ele não atua como roteador, e o tráfego roteado é negado de qualquer forma.
A saída do \texttt{iptables} confirma que a regra de liberação ficou restrita à interface
\texttt{eth1}, à origem 192.168.50.20, ao destino 192.168.50.10 e à porta 8080/TCP.

\subsection{Varredura depois do firewall}

Repetidas as varreduras, apenas a porta 8080 continua \emph{open}; as demais passam a
\emph{filtered}, pois os pacotes são descartados sem resposta.

\evidenciaextra{fase2_nmap_depois.txt}{\texttt{nmap} do cliente contra o servidor, depois do firewall}

\printtela{fase2-nmap-depois.png}{Execução manual no cliente: \texttt{nmap} depois do firewall, com somente 8080/tcp aberta e 65.534 portas filtradas (22/tcp passou de \emph{open} para \emph{filtered}); \texttt{curl} na 8080 responde HTTP 200 e \texttt{nc} na 22 expira}

% ===========================================================================
\section{Configuração de IAM e armazenamento em nuvem}\label{sec:iam}

\subsection{Plataforma utilizada}

Optou-se pela simulação local permitida pelo enunciado. O servidor e o cliente oficiais
do MinIO deixaram de ser distribuídos (o endereço oficial de \textit{download} responde
``410 Gone'') \cite{minio-arquivado}; por isso foi usado o SILO \cite{silo-docs}, versão
\texttt{RELEASE.2026-11-12}, instalado a partir dos pacotes oficiais do projeto com
verificação de SHA-256. O SILO mantém a API S3, a API de administração, as variáveis
\texttt{MINIO\_*} e os comandos do cliente (\texttt{mcli}), de modo que os conceitos e
comandos são os mesmos do MinIO. O serviço escuta na porta 9000 (API S3) e 9001 (console),
ambas bloqueadas para a rede pelo \textit{firewall} da \autoref{sec:firewall}; a
administração foi feita localmente no servidor.

As credenciais de administrador e do usuário de laboratório foram geradas aleatoriamente
na própria VM, gravadas em arquivos com permissão 600 e nunca aparecem nas evidências
(os comandos registram apenas o nome da variável, como \texttt{\$LAB\_SECRET}).

\subsection{Usuário e política de menor privilégio}

Foram criados dois \textit{buckets}: \texttt{lab-bucket}, ao qual o usuário terá acesso,
e \texttt{outro-bucket}, usado para provar que o acesso não se estende a outros recursos.
A política \texttt{lab-rw} concede somente o necessário para ``ler e escrever'' no
\texttt{lab-bucket} (\autoref{tab:politica}), aplicando o princípio do menor privilégio
\cite{nist80053}. Não há curingas de ação, não há permissão administrativa, e mesmo a
exclusão de objetos ficou de fora por não ser exigida pelo enunciado.

\begin{table}[htbp]
  \centering
  \caption{Permissões da política \texttt{lab-rw}}\label{tab:politica}
  \begin{tabular}{L{0.3\textwidth}L{0.36\textwidth}L{0.24\textwidth}}
    \hline
    \textbf{Ação} & \textbf{Recurso} & \textbf{Finalidade} \\
    \hline
    \texttt{s3:ListBucket}, \texttt{s3:GetBucketLocation} & \texttt{arn:aws:s3:::lab-bucket} & listar o conteúdo do \textit{bucket} \\
    \texttt{s3:GetObject} & \texttt{arn:aws:s3:::lab-bucket/*} & ler objetos \\
    \texttt{s3:PutObject} & \texttt{arn:aws:s3:::lab-bucket/*} & gravar objetos \\
    \hline
    \multicolumn{3}{L{0.9\textwidth}}{\footnotesize Não concedido: criar ou remover \textit{buckets}, apagar objetos, alterar políticas, acessar outros \textit{buckets} e qualquer ação \texttt{admin:*}.} \\
    \hline
  \end{tabular}
  \legend{Fonte: elaborado pelo autor.}
\end{table}

\printtela{fase3-silo-versao.png}{Execução manual no servidor: SILO e \texttt{mcli} instalados, versão \texttt{RELEASE.2026-11-12T00-00-00Z}}
\printtela{fase3-silo-servico.png}{Execução manual no servidor: senha gerada aleatoriamente em arquivo privado, SILO no ar (\texttt{health: HTTP 200}) e alias administrativo \texttt{root} configurado}

\printtela{fase3-buckets.png}{Execução manual no servidor: criação dos \textit{buckets} \texttt{lab-bucket} e \texttt{outro-bucket} e envio dos arquivos de teste}

\printtela{fase3-iam.png}{Execução manual no servidor: política \texttt{lab-rw} (somente \texttt{lab-bucket}), criação do \texttt{lab-user} e vínculo da política (\texttt{PolicyName: lab-rw})}

\evidenciaextra{fase3_iam.txt}{Instalação do SILO, criação dos buckets, da política \texttt{lab-rw} e do usuário \texttt{lab-user}}

\subsection{Bucket público e correção para privado}

Primeiro o \texttt{lab-bucket} foi deliberadamente tornado público com
\texttt{mcli anonymous set public}. Isso anexa ao \textit{bucket} uma política com
\texttt{"Principal": \{"AWS": ["*"]\}} e as ações de leitura, listagem, gravação e até
exclusão de objetos (\texttt{s3:DeleteObject}), ou seja, qualquer pessoa sem credenciais
pode listar, baixar, gravar e apagar objetos. As requisições \texttt{curl} sem assinatura
da evidência a seguir confirmam leitura, listagem e gravação anônimas.

\evidenciaextra{fase3_bucket_publico.txt}{\textit{Bucket} público: acesso anônimo de leitura, listagem e escrita}

\evidenciaresumo{manual_fase3_bucket_publico.txt}{Execução manual: \textit{bucket} público, com leitura, listagem e escrita anônimas (HTTP 200)}

Em seguida o \textit{bucket} foi corrigido com \texttt{mcli anonymous set private}, que
remove a política anônima. As mesmas requisições passam a receber \texttt{403 AccessDenied}.

\evidenciaextra{fase3_bucket_privado.txt}{\textit{Bucket} privado: acesso anônimo negado}

\evidenciaresumo{manual_fase3_bucket_privado.txt}{Execução manual: \textit{bucket} corrigido para privado, com as mesmas tentativas anônimas negadas (HTTP 403)}

A \autoref{tab:publico-privado} resume a diferença entre as duas configurações.

\begin{table}[htbp]
  \centering
  \caption{Diferença entre as configurações pública e privada}\label{tab:publico-privado}
  \begin{tabular}{L{0.3\textwidth}L{0.3\textwidth}L{0.3\textwidth}}
    \hline
    \textbf{Aspecto} & \textbf{Público} & \textbf{Privado} \\
    \hline
    Comando & \texttt{anonymous set public} & \texttt{anonymous set private} \\
    Política do \textit{bucket} (\texttt{get-json}) & \texttt{Allow} para \texttt{Principal "*"}: leitura, listagem, gravação e exclusão & \texttt{\{\}} (nenhuma política anônima) \\
    \texttt{GET} anônimo de objeto & HTTP 200 (conteúdo retornado) & HTTP 403 \texttt{AccessDenied} \\
    Listagem anônima & HTTP 200 (\texttt{ListBucketResult}) & HTTP 403 \texttt{AccessDenied} \\
    \texttt{PUT} anônimo & HTTP 200 (objeto gravado) & HTTP 403 \texttt{AccessDenied} \\
    \hline
  \end{tabular}
  \legend{Fonte: elaborado pelo autor, a partir das Evidências da \autoref{sec:iam}.}
\end{table}

% ===========================================================================
\section{Testes de verificação final}\label{sec:verificacao}

A verificação final repete os testes das fases anteriores com critérios explícitos de
aprovação: cada teste registra o resultado esperado e o veredito
(\texttt{PASSOU}/\texttt{FALHOU}), calculado pelo próprio \textit{script} a partir do
código de saída e do conteúdo da resposta.

\subsection{Portas acessíveis a partir do cliente}

\evidenciaresumo{fase4_cliente.txt}{Verificação a partir do cliente: somente a 8080 acessível}

\evidenciaextra{manual_fase4_cliente.txt}{Execução manual: varredura final com o SILO em execução no servidor, somente 8080/tcp aberta}

\printtela{fase4-cliente.png}{Execução manual no cliente: varredura final com 9000/9001 (SILO) \emph{filtered}, HTTP 200 na 8080, \emph{timeout} no SSH (22) e na API S3 (9000, código 28)}

\subsection{Bucket sem acesso público e usuário restrito}

Os testes de acesso anônimo e do usuário \texttt{lab-user} foram executados no próprio
servidor, contra a API local, já que a porta 9000 está bloqueada para a rede. Foram
testadas ações dentro do escopo da política (que devem funcionar) e seis ações fora dele
(que devem ser negadas).

\evidenciaresumo{fase4_servidor.txt}{Verificação no servidor: \textit{bucket} privado e \texttt{lab-user} limitado à política}

\printtela{fase4-servidor.png}{Execução manual no servidor: \textit{bucket} \texttt{private} (HTTP 403 anônimo); \texttt{lab-user} lista, lê e grava no \texttt{lab-bucket} e tem negadas as seis ações fora do escopo}

\chapter{Resultados e discussão}

A \autoref{tab:criterios} relaciona cada critério de avaliação do enunciado com o
resultado obtido e a evidência correspondente.

\begin{table}[htbp]
  \centering
  \caption{Critérios de avaliação e resultados}\label{tab:criterios}
  \begin{tabular}{L{0.3\textwidth}L{0.38\textwidth}L{0.2\textwidth}}
    \hline
    \textbf{Critério} & \textbf{Resultado} & \textbf{Evidência} \\
    \hline
    Cálculo de sub-redes e configuração de IP & 4 sub-redes /26 de 62 \textit{hosts}; VMs na sub-rede A; cliente na sub-rede B isolado do servidor & \autoref{sec:enderecamento} \\
    Efetividade do \textit{firewall} & antes: 22/tcp e 8080/tcp abertas; depois: somente 8080/tcp (65.534 portas filtradas) & \autoref{sec:firewall} \\
    Menor privilégio no IAM & \texttt{lab-user} lê, lista e grava só no \texttt{lab-bucket}; as 6 ações testadas fora do escopo foram negadas & \autoref{sec:iam}, \autoref{sec:verificacao} \\
    Armazenamento privado & acesso anônimo (leitura, listagem e escrita) passou de HTTP 200 para HTTP 403 & \autoref{sec:iam} \\
    Evidência completa da Fase 4 & 18 testes (13 no servidor e 5 no cliente), nenhuma falha & \autoref{sec:verificacao} \\
    \hline
  \end{tabular}
  \legend{Fonte: elaborado pelo autor.}
\end{table}

\section{Defesa em profundidade}

O laboratório mostra dois controles independentes protegendo o mesmo dado. Mesmo com o
\textit{bucket} público, a API S3 do SILO não era alcançável pelo cliente, porque o
\textit{firewall} só libera a porta 8080. Por outro lado, se a regra de \textit{firewall}
fosse relaxada, a política do \textit{bucket} e do usuário continuariam negando o acesso
indevido. Filtragem de rede responde \emph{quem alcança o serviço}; IAM responde
\emph{quem pode fazer o quê em qual recurso}. Um não substitui o outro.

\section{Equivalência com provedores de nuvem}

Os conceitos exercitados têm correspondência direta nos provedores públicos, o que
permite transpor o laboratório para uma conta real (\autoref{tab:nuvem})
\cite{nist800145,aws-s3-bpa}.

\begin{table}[htbp]
  \centering
  \caption{Correspondência entre o laboratório e a AWS}\label{tab:nuvem}
  \begin{tabular}{L{0.42\textwidth}L{0.48\textwidth}}
    \hline
    \textbf{Laboratório} & \textbf{AWS} \\
    \hline
    Sub-redes /26 em 192.168.50.0/24 & VPC com sub-redes \\
    \texttt{ufw} no servidor & \textit{Security Group} / \textit{Network ACL} \\
    Usuário \texttt{lab-user} + política \texttt{lab-rw} & Usuário ou \textit{role} do IAM + política gerenciada \\
    \texttt{mcli anonymous set public/private} & \textit{Bucket policy} com \texttt{Principal "*"} e \textit{S3 Block Public Access} \\
    \hline
  \end{tabular}
  \legend{Fonte: elaborado pelo autor.}
\end{table}

\section{Limitações}

\begin{alineas}
  \item a política de saída do \texttt{ufw} foi mantida como ``permitir'', pois o enunciado
    trata apenas do tráfego de entrada; em produção, a saída também deveria ser restrita;
  \item a exceção de gerência do Vagrant (porta 22 pela NAT, a partir do hospedeiro) existe
    apenas por causa da ferramenta de automação e não é alcançável pela rede do laboratório;
  \item os testes de acesso anônimo ao \textit{bucket} foram feitos a partir do próprio
    servidor, porque a rede bloqueia a porta 9000; isso isola o teste do controle de IAM
    do teste do controle de rede;
  \item o SILO é um \textit{fork} comunitário; em uma conta de nuvem real, o mesmo roteiro
    usaria IAM e S3 nativos do provedor.
\end{alineas}

\section{Observações da execução}

A varredura completa das 65.535 portas levou cerca de 26 segundos antes do
\textit{firewall} e cerca de 86 segundos depois. A diferença se explica pelo
comportamento de cada estado: uma porta \emph{fechada} responde imediatamente com um
pacote de recusa (\texttt{RST}), enquanto uma porta \emph{filtrada} não responde, e o
\texttt{nmap} precisa esperar o tempo limite de cada tentativa. Do ponto de vista de quem
varre, o \textit{firewall} portanto não só esconde os serviços como encarece o
reconhecimento da rede. O mesmo efeito aparece nos testes de aplicação: a conexão à
porta 22 expira por \textit{timeout}, em vez de ser recusada.

\chapter{Conclusão}

O laboratório atingiu os objetivos propostos. A rede 192.168.50.0/24 foi dividida em
sub-redes /26 e a ausência de comunicação entre sub-redes sem roteador foi demonstrada.
O \textit{firewall} do servidor, com política de negação por padrão e uma única regra
restrita por interface, origem, destino e porta, reduziu a superfície exposta ao cliente
de duas portas (22 e 8080) para apenas a 8080, como comprovado pelas varreduras com
\texttt{nmap}. No armazenamento, o \textit{bucket} configurado propositalmente como público
permitiu leitura, listagem e escrita sem autenticação; após a correção, todas essas
tentativas passaram a ser negadas. O usuário criado consegue apenas listar, ler e gravar
no \textit{bucket} autorizado, e todas as ações fora desse escopo foram negadas.

A principal lição é que controles de rede e de identidade se complementam: cada um cobre
a falha do outro, e a verificação com testes reproduzíveis, com resultado esperado
explícito, é o que permite afirmar que um controle realmente funciona.


% ----------------------------- Pós-textuais --------------------------------
\postextual
\bibliography{referencias}

% O enunciado pede "Anexos". Pela NBR 14724, material elaborado pelo próprio autor
% seria "Apêndice"; troque anexosenv por apendicesenv se o professor preferir.
\begin{anexosenv}
% Os arquivos de anexos/codigo/ são uma cópia do código usado na execução,
% feita automaticamente pelo rodar_lab.ps1 (ou: .\rodar_lab.ps1 -SomenteCopiarCodigo).
\newcommand{\codigo}[2]{%
  \par\medskip\noindent\textbf{#2} (\texttt{\detokenize{#1}})\par
  \IfFileExists{anexos/codigo/#1}{%
    \VerbatimInput[fontsize=\scriptsize,breaklines,breakanywhere,frame=single,framesep=2mm,numbers=left,numbersep=4pt]{anexos/codigo/#1}}{%
    \textcolor{red}{Arquivo anexos/codigo/\detokenize{#1} não encontrado.}}\par}

\chapter{Roteiro de execução e ambiente}\label{anexo:roteiro}

O laboratório completo é reproduzido pelos comandos abaixo, executados no hospedeiro
Windows a partir da pasta do projeto. Cada \textit{script} grava suas evidências em
\texttt{relatorio/evidencias/}.

\begin{Verbatim}[fontsize=\footnotesize,frame=single,framesep=2mm]
vagrant up              # cria e provisiona cliente e servidor
.\rodar_lab.ps1         # executa as fases 1 a 4 na ordem abaixo
\end{Verbatim}

\begin{table}[htbp]
  \centering
  \caption{Ordem de execução dos scripts}\label{tab:roteiro}
  \footnotesize
  \begin{tabular}{clll}
    \hline
    \textbf{Fase} & \textbf{VM} & \textbf{Comando} & \textbf{Evidência gerada} \\
    \hline
    1 & servidor & \texttt{sudo bash fase1\_ip.sh servidor 24}  & \texttt{fase1\_ip\_servidor\_24.txt} \\
    1 & cliente  & \texttt{sudo bash fase1\_ip.sh cliente 24}   & \texttt{fase1\_ip\_cliente\_24.txt} \\
    1 & cliente  & \texttt{bash fase1\_testes.sh 24}            & \texttt{fase1\_conectividade\_24.txt} \\
    1 & cliente  & \texttt{bash fase1\_subredes.sh}             & \texttt{fase1\_subredes.txt} \\
    1 & ambas    & \texttt{fase1\_ip.sh <papel> 26} e \texttt{fase1\_testes.sh 26} & \texttt{fase1\_*\_26.txt} \\
    1 & cliente  & \texttt{sudo bash fase1\_segmentacao.sh}     & \texttt{fase1\_segmentacao.txt} \\
    2 & servidor & \texttt{sudo bash fase2\_servicos.sh}        & \texttt{fase2\_servicos.txt} \\
    2 & cliente  & \texttt{sudo bash fase2\_scan.sh antes}      & \texttt{fase2\_nmap\_antes.txt} \\
    2 & servidor & \texttt{sudo bash fase2\_firewall.sh}        & \texttt{fase2\_firewall.txt} \\
    2 & cliente  & \texttt{sudo bash fase2\_scan.sh depois}     & \texttt{fase2\_nmap\_depois.txt} \\
    3 & servidor & \texttt{sudo bash fase3\_silo.sh}            & \texttt{fase3\_iam.txt} \\
    3 & servidor & \texttt{sudo bash fase3\_bucket.sh publico}  & \texttt{fase3\_bucket\_publico.txt} \\
    3 & servidor & \texttt{sudo bash fase3\_bucket.sh privado}  & \texttt{fase3\_bucket\_privado.txt} \\
    4 & servidor & \texttt{sudo bash fase4\_servidor.sh}        & \texttt{fase4\_servidor.txt} \\
    4 & cliente  & \texttt{sudo bash fase4\_cliente.sh}         & \texttt{fase4\_cliente.txt} \\
    \hline
  \end{tabular}
  \legend{Fonte: elaborado pelo autor.}
\end{table}

\ifcompleto\codigo{Vagrantfile}{Definição das VMs}\fi
\codigo{politica-lab-rw.json}{Política IAM de menor privilégio}

\chapter{Prints de configuração}
\renewcommand{\printw}{0.4}

\printtela{vagrant-up.png}{Recriação do ambiente: \texttt{vagrant destroy -f} e \texttt{vagrant up}, com o adaptador 1 em NAT e o adaptador 2 na rede interna (\texttt{intnet})}
\printtela{provisionamento.png}{Provisionamento concluído: instalação de \texttt{nmap}, \texttt{ufw}, \texttt{curl} e OpenSSH no servidor}
\printtela{acesso-servidor.png}{Acesso à VM servidor (\texttt{vagrant ssh servidor}): Ubuntu 22.04.5 LTS}
\printtela{acesso-cliente.png}{Acesso à VM cliente (\texttt{vagrant ssh cliente}): Ubuntu 22.04.5 LTS}
\printtela{virtualbox-rede-interna.png}{Adaptador 2 das VMs ligado à rede interna \texttt{lab-net}}
\printtela{vagrant-status.png}{VMs em execução (\texttt{vagrant status})}
\printtela{virtualbox-detalhes.png}{VirtualBox Manager: VM \texttt{lab-nsce-servidor} com o adaptador 1 em NAT e o adaptador 2 na rede interna \texttt{lab-net}}

\chapter{Scripts do laboratório}

\ifcompleto\else
Os \textit{scripts} que automatizam cada fase acompanham este relatório na pasta
\texttt{lab/} do projeto (versão completa com o código-fonte: \texttt{main-completo.pdf}).
A \autoref{tab:scripts} resume a função de cada um.

\begin{table}[htbp]
  \centering
  \caption{Scripts do laboratório}\label{tab:scripts}
  \footnotesize
  \begin{tabular}{L{0.32\textwidth}L{0.6\textwidth}}
    \hline
    \textbf{Script} & \textbf{Função} \\
    \hline
    \texttt{provision.sh} & instala pacotes e gera/autoriza a chave SSH \\
    \texttt{fase1\_ip.sh} & IP estático por \textit{netplan} (/24 ou /26) \\
    \texttt{fase1\_subredes.sh} & cálculo das sub-redes /26 \\
    \texttt{fase1\_testes.sh} & \texttt{ping} e \texttt{ssh} cliente $\rightarrow$ servidor \\
    \texttt{fase1\_segmentacao.sh} & cliente na sub-rede B sem alcançar o servidor \\
    \texttt{fase2\_servicos.sh} & serviço web na 8080 e SSH \\
    \texttt{fase2\_scan.sh} & \texttt{nmap} antes/depois do \textit{firewall} \\
    \texttt{fase2\_firewall.sh} & regras do \texttt{ufw} \\
    \texttt{fase3\_silo.sh} & instala o SILO, cria \textit{buckets}, política e usuário \\
    \texttt{fase3\_bucket.sh} & \textit{bucket} público e depois privado \\
    \texttt{fase4\_servidor.sh}, \texttt{fase4\_cliente.sh} & verificação final com resultado esperado e veredito \\
    \hline
  \end{tabular}
  \legend{Fonte: elaborado pelo autor.}
\end{table}
\fi

\ifcompleto
\codigo{comum.sh}{Funções comuns e registro de evidências}
\codigo{provision.sh}{Provisionamento das VMs}
\codigo{fase1_ip.sh}{Fase 1 --- IP estático}
\codigo{fase1_subredes.sh}{Fase 1 --- cálculo de sub-redes}
\codigo{fase1_testes.sh}{Fase 1 --- testes de conectividade}
\codigo{fase1_segmentacao.sh}{Fase 1 --- teste de segmentação}
\codigo{fase2_servicos.sh}{Fase 2 --- serviços}
\codigo{fase2_scan.sh}{Fase 2 --- varredura}
\codigo{fase2_firewall.sh}{Fase 2 --- firewall}
\codigo{fase3_silo.sh}{Fase 3 --- SILO e IAM}
\codigo{fase3_bucket.sh}{Fase 3 --- bucket público/privado}
\codigo{fase4_servidor.sh}{Fase 4 --- verificação no servidor}
\codigo{fase4_cliente.sh}{Fase 4 --- verificação no cliente}
\fi

\end{anexosenv}

\end{document}
